\documentclass[10pt,a4paper]{amsart}
\usepackage[margin=19mm]{geometry}
\usepackage{amsmath,amssymb,mathtools}
\usepackage[scr=rsfs,cal=euler]{mathalfa}
\usepackage{xfrac}
\usepackage[unicode,hidelinks,bookmarks=false]{hyperref}
\newcommand{\ie}{i.e.\ }
\newcommand{\cf}{cf.\ }
\newcommand{\resp}{respectively\ }
\newcommand{\Subsection}{\subsection}
\providecommand{\nobreakdash}{\nobreak\hbox{-}}
\setlength{\emergencystretch}{2em}
\multlinegap0pt
\usepackage{amssymb}
\usepackage{enumerate}
\usepackage{xcolor}
\newtheorem{lem}{Lemma}
\newcommand{\eps}{\varepsilon}
\title{Conjugacy co-amenability}
\author{Mehrdad Kalantar, Srivatsav Kunnawalkam Elayavalli}
\date{Source: arXiv:2602.14176v1 (CC BY 4.0)}
\begin{document}
\maketitle

\noindent\emph{Source and licence.} Conjugacy co-amenability. Version arXiv:2602.14176v1 (CC BY 4.0), distributed by arXiv under the Creative Commons Attribution 4.0 International licence, http://creativecommons.org/licenses/by/4.0/, as recorded in the arXiv licence metadata for this exact version and shown on the abstract page https://arxiv.org/abs/2602.14176v1. Full licence notice: \texttt{licenses/arxiv-2602.14176-CC-BY-4.0.txt}.\par\smallskip

\noindent\textit{Editorial scope.} Counted unit: the article's lem lem (source0.tex lines 313--315). The article's complete original proof of it is delivered with it (source0.tex lines 317--328, one \texttt{proof} environment of 12 lines). No mathematics is written, repaired or re-derived: the statement and the proof are the article's own lines, in the article's own notation, and no retained line is substituted or re-flowed.  No further source-local context is needed: every cross-reference of every retained line resolves inside the two delivered spans.  Nothing was omitted from any retained span, so the package carries no omission record.  No retained line contains a citation, so the wrapper carries no bibliography and no bibliography entry.  No retained line contains a figure environment, an image-inclusion command or drawing code, so no image is delivered and the package declares no asset dependency.  Editorial formatting only, in addition: an AMS \texttt{amsart} preamble that restates the article's own declarations and macros used by the retained lines, namely the environments \texttt{lem} on separate counters, together with the standard packages those lines need.  The retained environments print in this document as Lemma 1 in order of appearance, whereas the article prints the same items with the numbering of its own full document.  The complete native source of the article as distributed in the arXiv e-print for this exact version is bundled unchanged as \texttt{sources/194693/source0.tex}.
\begin{lem}
Let $G$ be a group and $H\le G$. If $G$ has a non-amenable subgroup $K\le G$ such that $K\cap gHg^{-1} = \{e\}$ for all $g\in G$, then $H$ is not co-amenable in $G$.
\end{lem}

\begin{proof}
Since $K$ is non-amenable, there exists a finite set $F\subset K$ and $\eps>0$ such that $\max_{t\in F} \|t\xi - \xi\|_2 \ge \eps\|\xi\|$ for every $\xi\in\ell^2(K)$. 

Let $S\subset G$ be a complete set of representatives for the double coset space $K/G\backslash H$ (that is, $G = \sqcup_{s\in S} KsH$). Note that for $k, k'\in K$ and $s, s'\in S$, if $ksH = k's'H$ then $s=s'$ and $s^{-1}k^{-1}k's\in H$, which implies $k=k'$ since $K\cap sHs^{-1} = \{e\}$.

Given $\eta\in\ell^2(G/H)$, for every $s\in S$ we define $\eta_s\in\ell^2(K)$ by $\eta_s(k):= \eta(ks)$ for every $k\in K$. 
 Then
\[
\max_{t\in F} \|t\eta - \eta\|_2^2 = \max_{t\in F} \sum_{s\in S}\|t\eta_s - \eta_s\|_2^2 \ge \eps^2\sum_{s\in S}\|\eta_s\|_2^2 =\eps^2\|\eta\|_2^2 .
\]
This shows that $H$ is not co-amenable in $G$.
\end{proof}

\end{document}
